\section{Related work}
\label{sec:related}

VegasCore sits at the meeting point of three literatures --- representations
of imperfect-information games, mechanized game theory, and the semantics and
verification of contract runtimes.  We position it against each in turn,
emphasizing what is shared and where the constructions genuinely diverge.

\paragraph{Imperfect-information game representations.}
The closest comparison objects are state-based and tree-based models of finite
games with hidden information.  Factored-observation stochastic games
(FOSGs)~\citep{Kovarik2022} present such games through per-player observation
functions over a shared world state; extensive-form games
\citep{vonNeumannMorgenstern1944,OsborneRubinstein1994} present them as trees
with information sets; and \TT{OpenSpiel}~\citep{LanctotEtAl2019OpenSpiel}
packages many such games behind a common solver-facing API.  VegasCore does not
compete with these as a game format: it \emph{denotes} a game from a program
and \emph{exports} both an FOSG and an EFG presentation, with payoff-faithful
adequacy theorems relating them to the native frontier game
(\Cref{thm:fosg,thm:efg}).  The divergence is on two axes.  On the information
axis VegasCore is strictly less general: FOSGs admit arbitrary per-player
observation factorings, whereas VegasCore fixes an owner-or-public visibility
discipline in which \(\Reveal\) is the only sealed-to-public channel, so a
private observation that is neither public nor owned by a single player is not
expressible (\Cref{sec:disc:limits}).  In exchange it adds what those
state-based and tree-based models leave out: an operational semantics, a
sealed-then-revealed temporal discipline enforced by the commit/reveal barrier
(\Cref{thm:barrier}), and a refinement order carrying the game down to a
concrete runtime.  Relative to \TT{OpenSpiel}, VegasCore is a
semantics-and-proof layer rather than a solver environment; its finite pure
kernel-game export (\Cref{sec:applications}) is a bridge \emph{toward} such
tools, not a replacement for them.  The graphical-model reading --- multi-agent
influence diagrams~\citep{koller2003multi} --- is incomparable for the reasons
set out in the probabilistic-programming remark of \Cref{sec:discussion}.

\paragraph{Mechanized game theory.}
Prior mechanizations have largely formalized the game model itself: existence of
Nash equilibrium in Coq and Isabelle~\citep{LeRoux2017NashCoqIsabelle} and, most
closely related, games of incomplete information built on belief functions in
the Coq proof assistant~\citep{PomeretCoquot2023BelGames}.  These efforts prove
theorems
\emph{about} games --- equilibrium existence, solution-concept properties ---
taking the game as the primitive object.  VegasCore's emphasis is orthogonal
and complementary: rather than a new existence theorem it provides a checked
\emph{path} from a program through compilation, scheduling, presentation, and
runtime descent, reusing standard finite-game existence results at the leaves
(\Cref{sec:applications}).  The contribution is the family of faithfulness
relations that carry a game-theoretic fact across those boundaries, not the
game theory proved at any one of them.

\paragraph{Event structures, refinement, and compositional structure.}
The event graph is a typed, finite, asymmetric variant of an event
structure~\citep{NielsenPlotkinWinskel1981,Winskel1986EventStructures}, and the
schedule-confluence results are the protocol-setting analogue of the classical
independence-permutation property.  The runtime layer is an instance of
forward-simulation refinement~\citep{LynchVaandrager1995,AbadiLamport1991},
specialized to stochastic transition kernels and extended with the
strategy-space side conditions needed to transport equilibria
(\Cref{sec:ref:cheap-talk}).  Compositional and categorical game
theory~\citep{GhaniHedgesWinschel2018CGT,Hedges2016CGT} approaches the same
preservation question from the opposite direction, building games by
composition; VegasCore's EU-preserving kernel-game morphisms suggest such
structure (\Cref{sec:disc:ext}), but here composition is a prospective
extension rather than the organizing principle.

\paragraph{Contract runtimes and EVM semantics.}
Contract-oriented domain-specific languages such as
Scilla~\citep{Sergey2019Scilla} and Marlowe~\citep{Lamela2020Marlowe} give
executable semantics to a restricted programming model, and KEVM provides a
complete executable semantics of the Ethereum Virtual Machine
itself~\citep{Hildenbrandt2018KEVM}.  VegasCore deliberately does not model any
of these: it talks about machines, observations, and stochastic transition
kernels above a generic refinement boundary, and such semantics are exactly the
concrete backends that would sit \emph{below} it
(\Cref{sec:refinement,sec:disc:related-artifacts}).  A backend package supplies
the state, observation, and transition model --- whether an EVM semantics, an
executable chain such as ConCertLean, or a payment-channel machine --- and
proves it projects onto the event-graph machine; VegasCore contributes the
game-preservation obligations that such a projection must discharge, not the
runtime model.  The frontrunning and miner-extractable-value
literature~\citep{Daian2020Flash} motivates the message-in-flight backend, whose
public pending-message buffers make exactly that pending-transaction visibility
strategically explicit (\Cref{sec:ref:runtime-examples}).
